The requested upper bound is **false**. In fact, the proposed extremizing “row-sign vectors” are not local maximizers of the \(\ell^r\)-operator ratio. Here is a rigorous perturbation showing that, for every integer \(r\ge 3\), one can get a ratio strictly larger than \(\Gamma_r\).

Fix \(i\in\{1,\dots,7\}\), and let

\[
c:=C e_i .
\]

Since \(C=C^\top\), this is also the \(i\)-th row-sign vector. Because \(C\mathbf 1=-\mathbf 1\), the vector \(c\) has sum \(-1\), hence exactly three entries equal to \(+1\) and four entries equal to \(-1\). Also,

\[
Cc=C^2 e_i=(8I_7-J_7)e_i=8e_i-\mathbf 1 .
\]

Thus the row-sign vector gives

\[
\frac{\|Cc\|_r^r}{\|c\|_r^r}
=
\frac{7^r+6}{7}
=
\Gamma_r^r.
\]

Now define a perturbation direction

\[
h:=\frac{7\mathbf 1+c}{24}.
\]

Equivalently,

\[
h_j=
\begin{cases}
\frac13, & c_j=1,\\[2mm]
\frac14, & c_j=-1.
\end{cases}
\]

Then

\[
c^\top h
=
3\cdot \frac13-4\cdot \frac14
=
0,
\]

so \(h\) is tangent to the \(\ell^r\)-sphere at \(c\) to first order. For \(0<\varepsilon<4\), set

\[
x_\varepsilon:=c+\varepsilon h.
\]

The denominator is

\[
\|x_\varepsilon\|_r^r
=
3\left(1+\frac{\varepsilon}{3}\right)^r
+
4\left(1-\frac{\varepsilon}{4}\right)^r .
\]

Next compute \(C h\). Using \(C\mathbf 1=-\mathbf 1\) and \(Cc=8e_i-\mathbf 1\),

\[
C h
=
\frac{7C\mathbf 1+Cc}{24}
=
\frac{-7\mathbf 1+8e_i-\mathbf 1}{24}
=
\frac{e_i-\mathbf 1}{3}.
\]

Therefore

\[
Cx_\varepsilon
=
Cc+\varepsilon C h
=
8e_i-\mathbf 1+\frac{\varepsilon}{3}(e_i-\mathbf 1).
\]

Its \(i\)-th coordinate is still \(7\), while the other six coordinates are

\[
-1-\frac{\varepsilon}{3}.
\]

Hence

\[
\|Cx_\varepsilon\|_r^r
=
7^r+6\left(1+\frac{\varepsilon}{3}\right)^r .
\]

Define

\[
R_r(\varepsilon)
:=
\frac{\|Cx_\varepsilon\|_r^r}{\|x_\varepsilon\|_r^r}
=
\frac{7^r+6\left(1+\frac{\varepsilon}{3}\right)^r}
{3\left(1+\frac{\varepsilon}{3}\right)^r
+
4\left(1-\frac{\varepsilon}{4}\right)^r}.
\]

At \(\varepsilon=0\),

\[
R_r(0)=\frac{7^r+6}{7}=\Gamma_r^r.
\]

But differentiating at \(\varepsilon=0\) gives

\[
R_r'(0)
=
\frac{2r}{7}
>
0.
\]

Therefore, for all sufficiently small \(\varepsilon>0\),

\[
R_r(\varepsilon)>\Gamma_r^r.
\]

Equivalently,

\[
\frac{\|Cx_\varepsilon\|_r}{\|x_\varepsilon\|_r}>\Gamma_r.
\]

Thus

\[
\|C\|_{r\to r}>\Gamma_r
\]

for every integer \(r\ge 3\). Consequently, the claimed sharp upper bound

\[
\|C\|_{r\to r}\le \Gamma_r
\]

cannot be true, and equality at the row-sign vectors cannot characterize the operator norm.